\documentclass[a4paper,11pt]{article}

\usepackage[a4paper,margin=2cm]{geometry}
\usepackage[T1]{fontenc}
\usepackage[utf8]{inputenc}
\usepackage[ngerman,english]{babel}
\usepackage{lmodern}
\usepackage{microtype}
\usepackage{xcolor}
\usepackage{parskip}
\usepackage{enumitem}
\usepackage[most]{tcolorbox}
\usepackage{titlesec}
\usepackage[hidelinks]{hyperref}
\usepackage{array}
\usepackage{longtable}
\usepackage[table]{xcolor}

\definecolor{HeaderBlueDark}{RGB}{70,110,170}
\definecolor{RowBlueLight}{RGB}{235,243,255}
\definecolor{RowBlueDark}{RGB}{220,233,250}
\definecolor{SoftGray}{RGB}{90,90,90}

\setlength{\parindent}{0pt}
\setlist[itemize]{leftmargin=1.4em,itemsep=0.35em,topsep=0.35em}
\linespread{1.06}
\renewcommand{\arraystretch}{1.35}
\setlength{\tabcolsep}{5pt}

\titleformat{\section}[block]
{\Large\bfseries\color{HeaderBlueDark}}
{}{0pt}{}
\titlespacing{\section}{0pt}{1.3em}{0.8em}

\titleformat{\subsection}[block]
{\large\bfseries\color{HeaderBlueDark}}
{}{0pt}{}
\titlespacing{\subsection}{0pt}{1.0em}{0.6em}

\newtcolorbox{exbox}{enhanced,breakable,colback=RowBlueLight,colframe=RowBlueDark,boxrule=0.4pt,arc=1.5mm,left=1.2mm,right=1.2mm,top=1mm,bottom=1mm,title={Beispiele \textnormal{(Examples)}},fonttitle=\bfseries,colbacktitle=RowBlueDark,coltitle=black}

\NewDocumentEnvironment{grammarentry}{mm+b}{%
    \begin{tcolorbox}[enhanced,breakable,colback=white,colframe=HeaderBlueDark,boxrule=0.6pt,arc=1.5mm,left=1.5mm,right=1.5mm,top=1mm,bottom=1mm,colbacktitle=HeaderBlueDark,coltitle=white,fonttitle=\bfseries,title={#1\hfill\normalfont\footnotesize #2}]
    #3
    \end{tcolorbox}
}{}

\newcommand{\GermanText}[1]{\textbf{Deutsch: }#1\par}
\newcommand{\EnglishText}[1]{{\small\color{SoftGray}\textbf{English: }#1\par}}
\newcommand{\ExampleItem}[2]{\item \textbf{#1}\\{\small\color{SoftGray}#2}}

\title{Die n-Deklination \textnormal{(The n-declension)}}
\date{}

\begin{document}
\renewcommand{\contentsname}{Inhaltsverzeichnis \textnormal{(Table of Contents)}}
\maketitle
\tableofcontents
\newpage

\section{Grundidee \textnormal{(Basic idea)}}

\begin{grammarentry}{Was ist die n-Deklination? \textnormal{(What is the n-declension?)}}{schwache Deklination maskuliner Nomen \textnormal{(weak declension of masculine nouns)}}
\GermanText{Die \textbf{n-Deklination} ist eine besondere Deklination bestimmter \textbf{maskuliner Nomen}. Diese Nomen bekommen im Singular in \textbf{Akkusativ, Dativ und Genitiv} normalerweise die Endung \textbf{-n} oder \textbf{-en}. Nur der \textbf{Nominativ Singular} bleibt ohne diese zusätzliche Endung.}
\EnglishText{The \textbf{n-declension} is a special declension of certain \textbf{masculine nouns}. In the singular, these nouns normally take the ending \textbf{-n} or \textbf{-en} in the \textbf{accusative, dative, and genitive}. Only the \textbf{nominative singular} remains without this additional ending.}

\GermanText{Ein typisches Beispiel ist \textbf{der Student}: der Student, den Studenten, dem Studenten, des Studenten.}
\EnglishText{A typical example is \textbf{der Student}: der Student, den Studenten, dem Studenten, des Studenten.}
\end{grammarentry}

\begin{exbox}
\begin{itemize}
\ExampleItem{Der Student kommt.}{The student is coming.}
\ExampleItem{Ich sehe den Studenten.}{I see the student.}
\ExampleItem{Ich helfe dem Studenten.}{I help the student.}
\ExampleItem{Das Buch des Studenten liegt hier.}{The student's book is here.}
\end{itemize}
\end{exbox}

\section{Die wichtigste Regel \textnormal{(The most important rule)}}

\begin{grammarentry}{Singular \textnormal{(Singular)}}{Nominativ anders, drei andere Kasus mit -(e)n \textnormal{(nominative differs; three other cases take -(e)n)}}
\GermanText{Merksatz: \textbf{Nominativ Singular ohne -(e)n; Akkusativ, Dativ und Genitiv Singular mit -(e)n.}}
\EnglishText{Memory rule: \textbf{Nominative singular without -(e)n; accusative, dative, and genitive singular with -(e)n.}}

\GermanText{Die Endung gehört zum \textbf{Nomen selbst}. Sie ist also nicht die Endung des Artikels oder eines Adjektivs.}
\EnglishText{The ending belongs to the \textbf{noun itself}. It is therefore not the ending of the article or an adjective.}
\end{grammarentry}

\rowcolors{2}{RowBlueLight}{RowBlueDark}
\begin{longtable}{>{\bfseries}p{2.8cm}|p{4.4cm}|p{7.0cm}}
\rowcolor{HeaderBlueDark}
\color{white}\textbf{Kasus (Case)} & \color{white}\textbf{Form (Form)} & \color{white}\textbf{Beispiel / Übersetzung (Example / Translation)} \\
Nominativ & der Student & Der Student lernt. / The student studies. \\
Akkusativ & den Studenten & Ich sehe den Studenten. / I see the student. \\
Dativ & dem Studenten & Ich helfe dem Studenten. / I help the student. \\
Genitiv & des Studenten & Das Buch des Studenten. / The student's book. \\
\end{longtable}

\section{Wann benutzt man -n und wann -en? \textnormal{(When do we use -n and when -en?)}}

\begin{grammarentry}{Endung \textnormal{(Ending)}}{Form des Wortes \textnormal{(form of the word)}}
\GermanText{Viele n-Deklination-Nomen, die im Nominativ Singular bereits auf \textbf{-e} enden, bekommen normalerweise nur \textbf{-n}: der Junge $\rightarrow$ den Jungen; der Kollege $\rightarrow$ dem Kollegen; der Löwe $\rightarrow$ des Löwen.}
\EnglishText{Many n-declension nouns that already end in \textbf{-e} in the nominative singular normally add only \textbf{-n}: der Junge $\rightarrow$ den Jungen; der Kollege $\rightarrow$ dem Kollegen; der Löwe $\rightarrow$ des Löwen.}

\GermanText{Viele andere schwache Maskulina bekommen \textbf{-en}: der Student $\rightarrow$ den Studenten; der Journalist $\rightarrow$ dem Journalisten; der Elefant $\rightarrow$ des Elefanten.}
\EnglishText{Many other weak masculine nouns take \textbf{-en}: der Student $\rightarrow$ den Studenten; der Journalist $\rightarrow$ dem Journalisten; der Elefant $\rightarrow$ des Elefanten.}

\GermanText{Für Lernende ist es am sichersten, das Nomen zusammen mit einer gebeugten Form zu lernen, zum Beispiel: \textbf{der Student -- des Studenten -- die Studenten}.}
\EnglishText{For learners, the safest method is to learn the noun together with an inflected form, for example: \textbf{der Student -- des Studenten -- die Studenten}.}
\end{grammarentry}

\section{Welche Nomen gehören häufig zur n-Deklination? \textnormal{(Which nouns often belong to the n-declension?)}}

\subsection{Maskuline Personen- und Tierbezeichnungen auf -e \textnormal{(Masculine words for people and animals ending in -e)}}

\begin{grammarentry}{Typische Gruppe \textnormal{(Typical group)}}{Endung -e \textnormal{(ending -e)}}
\GermanText{Viele maskuline Personen- und Tierbezeichnungen auf \textbf{-e} gehören zur n-Deklination.}
\EnglishText{Many masculine nouns referring to people and animals and ending in \textbf{-e} belong to the n-declension.}
\end{grammarentry}

\rowcolors{2}{RowBlueLight}{RowBlueDark}
\begin{longtable}{>{\bfseries}p{4.0cm}|p{5.0cm}|p{5.0cm}}
\rowcolor{HeaderBlueDark}
\color{white}\textbf{Nominativ (Nominative)} & \color{white}\textbf{Andere Kasus (Other cases)} & \color{white}\textbf{English} \\
der Junge & den/dem/des Jungen & boy \\
der Kollege & den/dem/des Kollegen & colleague \\
der Kunde & den/dem/des Kunden & customer \\
der Experte & den/dem/des Experten & expert \\
der Löwe & den/dem/des Löwen & lion \\
der Hase & den/dem/des Hasen & hare / rabbit \\
der Zeuge & den/dem/des Zeugen & witness \\
der Neffe & den/dem/des Neffen & nephew \\
\end{longtable}

\subsection{Fremdwörter mit typischen Endungen \textnormal{(Loanwords with typical endings)}}

\begin{grammarentry}{Sehr häufige Muster \textnormal{(Very common patterns)}}{-ant, -ent, -ist und weitere \textnormal{(-ant, -ent, -ist and others)}}
\GermanText{Viele maskuline Fremdwörter, besonders Personenbezeichnungen, gehören zur n-Deklination. Häufige Endungen sind zum Beispiel \textbf{-ant, -ent, -ist, -oge, -at, -et, -nom}.}
\EnglishText{Many masculine loanwords, especially nouns referring to people, belong to the n-declension. Common endings include \textbf{-ant, -ent, -ist, -oge, -at, -et, -nom}.}
\end{grammarentry}

\rowcolors{2}{RowBlueLight}{RowBlueDark}
\begin{longtable}{>{\bfseries}p{3.5cm}|p{4.5cm}|p{6.0cm}}
\rowcolor{HeaderBlueDark}
\color{white}\textbf{Muster (Pattern)} & \color{white}\textbf{Beispiel (Example)} & \color{white}\textbf{Form außerhalb Nom. Sg. (Form outside nom. sg.)} \\
-ant & der Praktikant & den/dem/des Praktikanten \\
-ent & der Student & den/dem/des Studenten \\
-ist & der Journalist & den/dem/des Journalisten \\
-oge & der Biologe & den/dem/des Biologen \\
-at & der Diplomat & den/dem/des Diplomaten \\
-et & der Athlet & den/dem/des Athleten \\
-nom & der Astronom & den/dem/des Astronomen \\
\end{longtable}

\subsection{Wichtige Wörter ohne offensichtliche Endung \textnormal{(Important words without an obvious ending)}}

\begin{grammarentry}{Auswendig lernen \textnormal{(Learn by heart)}}{häufige Einzelwörter \textnormal{(common individual nouns)}}
\GermanText{Einige sehr häufige schwache Maskulina erkennt man nicht zuverlässig an einer typischen Endung. Diese sollte man als Wortschatz lernen.}
\EnglishText{Some very common weak masculine nouns cannot be reliably recognized by a typical ending. These should be learned as vocabulary.}
\end{grammarentry}

\rowcolors{2}{RowBlueLight}{RowBlueDark}
\begin{longtable}{>{\bfseries}p{4.0cm}|p{5.0cm}|p{5.0cm}}
\rowcolor{HeaderBlueDark}
\color{white}\textbf{Nominativ (Nominative)} & \color{white}\textbf{Andere Kasus (Other cases)} & \color{white}\textbf{English} \\
der Mensch & den/dem/des Menschen & human being / person \\
der Herr & den/dem/des Herrn & gentleman / Mr. \\
der Bär & den/dem/des Bären & bear \\
der Held & den/dem/des Helden & hero \\
der Prinz & den/dem/des Prinzen & prince \\
der Graf & den/dem/des Grafen & count \\
\end{longtable}

\section{Der Plural \textnormal{(The plural)}}

\begin{grammarentry}{Pluralformen \textnormal{(Plural forms)}}{nicht mit der Singularregel verwechseln \textnormal{(do not confuse them with the singular rule)}}
\GermanText{Bei vielen n-Deklination-Nomen endet auch der Plural auf \textbf{-n} oder \textbf{-en}. Deshalb sehen viele Pluralformen genauso aus wie Akkusativ, Dativ und Genitiv Singular.}
\EnglishText{For many n-declension nouns, the plural also ends in \textbf{-n} or \textbf{-en}. Therefore, many plural forms look the same as the accusative, dative, and genitive singular.}

\GermanText{Beispiel: \textbf{der Student} -- Singular: der Student / den Studenten / dem Studenten / des Studenten; Plural: die Studenten / die Studenten / den Studenten / der Studenten.}
\EnglishText{Example: \textbf{der Student} -- singular: der Student / den Studenten / dem Studenten / des Studenten; plural: die Studenten / die Studenten / den Studenten / der Studenten.}
\end{grammarentry}

\section{Artikel und Adjektive ändern die n-Deklination nicht \textnormal{(Articles and adjectives do not cancel the n-declension)}}

\begin{grammarentry}{Mehrere Endungen gleichzeitig \textnormal{(Several endings at the same time)}}{Artikel, Adjektiv und Nomen \textnormal{(article, adjective and noun)}}
\GermanText{Artikel und Adjektive werden nach ihren eigenen Regeln dekliniert. Das n-Deklination-Nomen bekommt trotzdem seine Endung.}
\EnglishText{Articles and adjectives are declined according to their own rules. The n-declension noun still receives its own ending.}
\end{grammarentry}

\begin{exbox}
\begin{itemize}
\ExampleItem{Ich sehe den jungen Studenten.}{I see the young student.}
\ExampleItem{Ich spreche mit einem netten Kollegen.}{I am speaking with a nice colleague.}
\ExampleItem{Das ist das Auto meines neuen Nachbarn.}{This is my new neighbor's car.}
\end{itemize}
\end{exbox}

\begin{grammarentry}{Achtung bei \glqq Nachbar\grqq \textnormal{(Note on \glqq Nachbar\grqq)}}{Gebrauchsvariation \textnormal{(variation in usage)}}
\GermanText{Bei einzelnen Wörtern können Wörterbücher oder Varietäten verschiedene zulässige Formen angeben. Wenn du unsicher bist, prüfe die Wörterbuchangabe für genau dieses Nomen. Die n-Deklination ist eine lexikalische Eigenschaft: Man kann sie nicht bei jedem maskulinen Nomen nur aus der Form erraten.}
\EnglishText{For some individual nouns, dictionaries or language varieties may list more than one acceptable form. If you are unsure, check the dictionary entry for that exact noun. The n-declension is a lexical property: it cannot be predicted from form alone for every masculine noun.}
\end{grammarentry}

\section{Besonders wichtig: \glqq Herr\grqq \textnormal{(Especially important: \glqq Herr\grqq)}}

\begin{grammarentry}{der Herr $\rightarrow$ Herrn \textnormal{(der Herr $\rightarrow$ Herrn)}}{sehr häufig im Alltag \textnormal{(very common in everyday German)}}
\GermanText{\textbf{Herr} hat die Formen: Nominativ \textbf{der Herr}, aber Akkusativ \textbf{den Herrn}, Dativ \textbf{dem Herrn}, Genitiv \textbf{des Herrn}.}
\EnglishText{\textbf{Herr} has the forms: nominative \textbf{der Herr}, but accusative \textbf{den Herrn}, dative \textbf{dem Herrn}, genitive \textbf{des Herrn}.}
\end{grammarentry}

\begin{exbox}
\begin{itemize}
\ExampleItem{Herr Müller kommt um zehn Uhr.}{Mr. Müller is coming at ten o'clock.}
\ExampleItem{Ich sehe Herrn Müller.}{I see Mr. Müller.}
\ExampleItem{Ich spreche mit Herrn Müller.}{I am speaking with Mr. Müller.}
\ExampleItem{Das Büro Herrn Müllers ist dort.}{Mr. Müller's office is there.}
\end{itemize}
\end{exbox}

\section{Gemischte Deklination: ähnliche, aber nicht identische Wörter \textnormal{(Mixed declension: similar but not identical nouns)}}

\begin{grammarentry}{Nicht reine n-Deklination \textnormal{(Not pure n-declension)}}{Genitiv auf -ns \textnormal{(genitive ending in -ns)}}
\GermanText{Einige maskuline Nomen sehen fast wie n-Deklination aus, haben aber im Genitiv Singular zusätzlich ein \textbf{-s}: zum Beispiel \textbf{der Name -- den Namen -- dem Namen -- des Namens}. Diese Gruppe nennt man oft \textbf{gemischte Deklination}.}
\EnglishText{Some masculine nouns look almost like n-declension nouns, but add an extra \textbf{-s} in the genitive singular: for example \textbf{der Name -- den Namen -- dem Namen -- des Namens}. This group is often called \textbf{mixed declension}.}
\end{grammarentry}

\rowcolors{2}{RowBlueLight}{RowBlueDark}
\begin{longtable}{>{\bfseries}p{3.7cm}|p{3.7cm}|p{3.7cm}|p{3.7cm}}
\rowcolor{HeaderBlueDark}
\color{white}\textbf{Nominativ (Nom.)} & \color{white}\textbf{Akkusativ (Acc.)} & \color{white}\textbf{Dativ (Dat.)} & \color{white}\textbf{Genitiv (Gen.)} \\
der Name & den Namen & dem Namen & des Namens \\
der Gedanke & den Gedanken & dem Gedanken & des Gedankens \\
der Glaube & den Glauben & dem Glauben & des Glaubens \\
der Wille & den Willen & dem Willen & des Willens \\
der Buchstabe & den Buchstaben & dem Buchstaben & des Buchstabens \\
\end{longtable}

\begin{grammarentry}{Merke \textnormal{(Remember)}}{Unterschied im Genitiv \textnormal{(difference in the genitive)}}
\GermanText{Reine n-Deklination: \textbf{des Studenten}. Gemischte Deklination: \textbf{des Namens}. Der wichtigste sichtbare Unterschied ist hier das zusätzliche \textbf{-s} im Genitiv.}
\EnglishText{Pure n-declension: \textbf{des Studenten}. Mixed declension: \textbf{des Namens}. The main visible difference here is the additional \textbf{-s} in the genitive.}
\end{grammarentry}

\section{Nicht mit substantivierten Adjektiven verwechseln \textnormal{(Do not confuse it with nominalized adjectives)}}

\begin{grammarentry}{der Deutsche, ein Deutscher \textnormal{(the German, a German)}}{Adjektivdeklination statt n-Deklination \textnormal{(adjective declension instead of n-declension)}}
\GermanText{Wörter wie \textbf{der Deutsche, der Erwachsene, der Bekannte, der Angestellte} sehen teilweise ähnlich aus, folgen aber der \textbf{Adjektivdeklination}. Ihre Endung hängt deshalb vom Artikel und Kasus ab.}
\EnglishText{Words such as \textbf{der Deutsche, der Erwachsene, der Bekannte, der Angestellte} may look similar, but they follow \textbf{adjective declension}. Their ending therefore depends on the article and case.}
\end{grammarentry}

\begin{exbox}
\begin{itemize}
\ExampleItem{Der Deutsche spricht Deutsch.}{The German speaks German.}
\ExampleItem{Ein Deutscher spricht Deutsch.}{A German speaks German.}
\ExampleItem{Ich sehe den Deutschen.}{I see the German.}
\ExampleItem{Ich sehe einen Deutschen.}{I see a German.}
\end{itemize}
\end{exbox}

\section{Auch nicht verwechseln: \glqq das Herz\grqq \textnormal{(Also do not confuse: \glqq das Herz\grqq)}}

\begin{grammarentry}{Sonderfall \textnormal{(Special case)}}{Neutrum, keine n-Deklination \textnormal{(neuter, not n-declension)}}
\GermanText{\textbf{das Herz} ist neutral und deshalb kein normales n-Deklination-Nomen. Es hat besondere Formen: \textbf{das Herz, des Herzens, dem Herzen, das Herz}.}
\EnglishText{\textbf{das Herz} is neuter and therefore not a normal n-declension noun. It has special forms: \textbf{das Herz, des Herzens, dem Herzen, das Herz}.}
\end{grammarentry}

\section{Wie erkenne ich ein n-Deklination-Nomen? \textnormal{(How do I recognize an n-declension noun?)}}

\begin{grammarentry}{Praktische Methode \textnormal{(Practical method)}}{Schritt für Schritt \textnormal{(step by step)}}
\GermanText{Benutze diese Reihenfolge:}
\EnglishText{Use this sequence:}
\begin{itemize}
\item \textbf{1.} Ist das Nomen maskulin? Wenn nein, ist es normalerweise keine n-Deklination.\\{\small\color{SoftGray}Is the noun masculine? If not, it is normally not an n-declension noun.}
\item \textbf{2.} Bezeichnet es eine Person oder ein Tier und endet es auf -e? Dann ist n-Deklination sehr wahrscheinlich.\\{\small\color{SoftGray}Does it refer to a person or animal and end in -e? Then n-declension is very likely.}
\item \textbf{3.} Hat es eine typische Fremdwortendung wie -ant, -ent, -ist, -oge, -at, -et oder -nom? Dann prüfe n-Deklination.\\{\small\color{SoftGray}Does it have a typical loanword ending such as -ant, -ent, -ist, -oge, -at, -et, or -nom? Then check for n-declension.}
\item \textbf{4.} Bei häufigen Einzelwörtern wie Mensch, Herr, Held, Prinz und Bär: Form mitlernen.\\{\small\color{SoftGray}For common individual nouns such as Mensch, Herr, Held, Prinz, and Bär: learn the inflected form together with the word.}
\item \textbf{5.} Im Zweifel Wörterbuch prüfen. Es gibt keine einzige Endungsregel, die alle Fälle sicher erkennt.\\{\small\color{SoftGray}If in doubt, check a dictionary. There is no single ending rule that identifies every case with certainty.}
\end{itemize}
\end{grammarentry}

\section{Typische Fehler \textnormal{(Common mistakes)}}

\begin{grammarentry}{Fehler 1 \textnormal{(Mistake 1)}}{Endung im Akkusativ vergessen \textnormal{(forgetting the accusative ending)}}
\GermanText{Falsch: \textbf{Ich sehe den Student.} Richtig: \textbf{Ich sehe den Studenten.}}
\EnglishText{Wrong: \textbf{Ich sehe den Student.} Correct: \textbf{Ich sehe den Studenten.}}
\end{grammarentry}

\begin{grammarentry}{Fehler 2 \textnormal{(Mistake 2)}}{Endung im Dativ vergessen \textnormal{(forgetting the dative ending)}}
\GermanText{Falsch: \textbf{Ich spreche mit dem Kollege.} Richtig: \textbf{Ich spreche mit dem Kollegen.}}
\EnglishText{Wrong: \textbf{Ich spreche mit dem Kollege.} Correct: \textbf{Ich spreche mit dem Kollegen.}}
\end{grammarentry}

\begin{grammarentry}{Fehler 3 \textnormal{(Mistake 3)}}{Endung im Genitiv vergessen \textnormal{(forgetting the genitive ending)}}
\GermanText{Falsch: \textbf{Das Auto des Student.} Richtig: \textbf{Das Auto des Studenten.}}
\EnglishText{Wrong: \textbf{Das Auto des Student.} Correct: \textbf{Das Auto des Studenten.}}
\end{grammarentry}

\begin{grammarentry}{Fehler 4 \textnormal{(Mistake 4)}}{-en im Nominativ Singular benutzen \textnormal{(using -en in nominative singular)}}
\GermanText{Falsch: \textbf{Der Studenten kommt.} Richtig: \textbf{Der Student kommt.}}
\EnglishText{Wrong: \textbf{Der Studenten kommt.} Correct: \textbf{Der Student kommt.}}
\end{grammarentry}

\begin{grammarentry}{Fehler 5 \textnormal{(Mistake 5)}}{Jedes maskuline Nomen als n-Deklination behandeln \textnormal{(treating every masculine noun as n-declension)}}
\GermanText{Nicht jedes maskuline Nomen ist schwach. Zum Beispiel: \textbf{der Tisch -- den Tisch -- dem Tisch -- des Tisches}. Hier gibt es keine n-Deklination.}
\EnglishText{Not every masculine noun is weak. For example: \textbf{der Tisch -- den Tisch -- dem Tisch -- des Tisches}. There is no n-declension here.}
\end{grammarentry}

\section{Präpositionen machen die n-Endung sichtbar \textnormal{(Prepositions often make the n-ending visible)}}

\begin{grammarentry}{Kasus nach Präpositionen \textnormal{(Case after prepositions)}}{n-Deklination bleibt aktiv \textnormal{(n-declension still applies)}}
\GermanText{Wenn eine Präposition Akkusativ oder Dativ verlangt, muss ein n-Deklination-Nomen ebenfalls seine Endung bekommen.}
\EnglishText{When a preposition requires the accusative or dative, an n-declension noun must also receive its ending.}
\end{grammarentry}

\begin{exbox}
\begin{itemize}
\ExampleItem{für den Studenten}{for the student}
\ExampleItem{ohne den Kollegen}{without the colleague}
\ExampleItem{mit dem Journalisten}{with the journalist}
\ExampleItem{von dem / vom Kunden}{from the customer}
\ExampleItem{bei dem / beim Experten}{with / at the expert}
\end{itemize}
\end{exbox}

\section{Relativsätze und Pronomen ändern die Regel nicht \textnormal{(Relative clauses and pronouns do not change the rule)}}

\begin{grammarentry}{Das Nomen bleibt schwach \textnormal{(The noun remains weak)}}{Kasus entscheidet die Form \textnormal{(case determines the form)}}
\GermanText{Auch in komplexeren Sätzen entscheidet der Kasus über die Form des schwachen Nomens.}
\EnglishText{Even in more complex sentences, the case determines the form of the weak noun.}
\end{grammarentry}

\begin{exbox}
\begin{itemize}
\ExampleItem{Der Student, den ich kenne, wohnt in Wien.}{The student whom I know lives in Vienna.}
\ExampleItem{Der Kollege, mit dem ich arbeite, ist freundlich.}{The colleague with whom I work is friendly.}
\ExampleItem{Der Journalist, dessen Artikel ich lese, arbeitet hier.}{The journalist whose article I read works here.}
\end{itemize}
\end{exbox}

\section{Lernstrategie \textnormal{(Learning strategy)}}

\begin{grammarentry}{So lernst du die Wörter \textnormal{(How to learn the nouns)}}{nicht nur den Nominativ lernen \textnormal{(do not learn only the nominative)}}
\GermanText{Lerne ein neues schwaches Nomen nicht nur als \textbf{der Student}. Lerne besser drei Informationen zusammen: \textbf{der Student -- des Studenten -- die Studenten}. Damit erkennst du sofort Genitiv und Plural und weißt, dass das Wort zur n-Deklination gehört.}
\EnglishText{Do not learn a new weak noun only as \textbf{der Student}. It is better to learn three pieces of information together: \textbf{der Student -- des Studenten -- die Studenten}. This immediately shows you the genitive and plural and tells you that the word belongs to the n-declension.}
\end{grammarentry}

\begin{grammarentry}{Minimalregel für B1 \textnormal{(Minimum rule for B1)}}{das musst du sicher können \textnormal{(what you should know confidently)}}
\GermanText{Für B1 solltest du sicher erkennen und bilden können: \textbf{der Student -- den Studenten -- dem Studenten -- des Studenten}; außerdem häufige Wörter wie \textbf{Junge, Kollege, Kunde, Mensch, Herr, Journalist, Student, Praktikant, Experte}.}
\EnglishText{For B1, you should be able to recognize and form with confidence: \textbf{der Student -- den Studenten -- dem Studenten -- des Studenten}; also common nouns such as \textbf{Junge, Kollege, Kunde, Mensch, Herr, Journalist, Student, Praktikant, Experte}.}
\end{grammarentry}

\section{Zusammenfassung \textnormal{(Summary)}}

\begin{grammarentry}{Merksatz \textnormal{(Memory rule)}}{kurz und wichtig \textnormal{(short and important)}}
\GermanText{\textbf{n-Deklination = bestimmte maskuline Nomen. Im Singular: nur Nominativ ohne -(e)n; Akkusativ, Dativ und Genitiv mit -(e)n.}}
\EnglishText{\textbf{n-declension = certain masculine nouns. In the singular: only the nominative has no -(e)n; accusative, dative, and genitive take -(e)n.}}

\GermanText{Beispiel: \textbf{der Student -- den Studenten -- dem Studenten -- des Studenten}.}
\EnglishText{Example: \textbf{der Student -- den Studenten -- dem Studenten -- des Studenten}.}

\GermanText{Sondergruppe: \textbf{der Name -- den Namen -- dem Namen -- des Namens} gehört zur gemischten Deklination, nicht zur reinen n-Deklination.}
\EnglishText{Special group: \textbf{der Name -- den Namen -- dem Namen -- des Namens} belongs to mixed declension, not pure n-declension.}
\end{grammarentry}

\end{document}
